\section{Discussion et conclusion}\label{sec:discussion}

\subsection*{Ce que les tensions révèlent sur les deux positions}

La mise en perspective des quatre tensions épistémologiques ci-dessus confirme que ni la position interdictionniste ni la position intégrationniste n'y apportent de réponse adéquate.

L'interdiction échoue à traiter les tensions 1 et 4: elle préserve la forme des évaluations existantes sans interroger leur validité. 
Elle suppose que nos critères actuels sont sains, et que le problème vient uniquement de l'outil. 
Mais si une machine peut satisfaire ces critères, l'outil n'est pas le problème: il est le révélateur d'un problème antérieur. 
De plus, elle est en contradiction interne avec la définition même de la compétence: si être compétent consiste à mobiliser des ressources pertinentes, interdire l'accès à une ressource pertinente est une décision pédagogique qui exige une justification --- pas un simple décret.

L'intégration non réfléchie échoue à traiter les tensions 2 et 3: sans redéfinition des critères et sans rendre le processus d'apprentissage observable, elle risque de certifier des performances sans compétences réelles, et d'abandonner les apprenants dans un écosystème informationnel où l'autorité du savoir est dispersée et opaque. 
Certifier que quelqu'un sait utiliser une IAG n'est pas la même chose que certifier qu'il a développé une pensée critique: la confusion entre les deux serait une régression, pas un progrès.

\subsection*{Prendre position: ce que le système scolaire doit affronter}

Le défi posé par les IAG est avant tout \emph{pédagogique et épistémologique}, non technologique. 
Il contraint le système scolaire à répondre à des questions qu'il avait pu différer jusqu'ici: que voulons-nous réellement que les apprenants développent?
Comment évaluer cela de façon valide? 
Quel est le rôle de l'école dans une société où la connaissance est de plus en plus externalisée dans des outils?

\cite{alexandre2019ia} formule cette question de façon provocatrice: dans une économie de la connaissance, l'école \og{}en tant que technologie de transmission de l'intelligence est d'ores et déjà une technologie dépassée\fg{}. %chktex 2
Sa conclusion --- faire de l'éducation une \og{}branche de la médecine\fg{}, centrée sur l'augmentation cognitive individuelle --- est discutable dans ses modalités. 
Mais le diagnostic mérite d'être pris au sérieux plutôt que mis de côté: si l'école se définit principalement par la transmission de savoirs formalisés et la certification de leur restitution, alors effectivement les IAG rendent cette définition obsolète. 
La réponse n'est pas de devenir une branche de la médecine, mais de déplacer l'accent de la \emph{restitution} vers la \emph{production}: former des apprenants capables non seulement de mobiliser des savoirs existants, mais de les questionner, de les situer, et d'en produire de nouveaux --- des capacités que les IAG amplifient mais ne remplacent pas.

La complexité du phénomène tient aussi au fait que les IAG ne forment pas un bloc homogène. 
Leurs effets varient considérablement selon les niveaux scolaires, les disciplines et les types de tâches: une IAG qui aide un étudiant de master à structurer une argumentation complexe n'a pas le même effet formateur --- ni le même risque --- qu'une IAG utilisée à la place d'un exercice de résolution de problèmes en première année. 
Une réponse pédagogique nuancée doit tenir compte de cette diversité plutôt que de chercher une politique universelle.

\subsection*{Réponse à la question posée, et ouverture}

Que deviennent les notions de compétence et de réussite scolaire dans un monde où les IAG sont en plein essor? 
Notre analyse permet de formuler une réponse précise: ces notions ne disparaissent pas, mais elles \emph{se complexifient} en révélant des tensions que le système scolaire portait déjà, et que la technologie contraint désormais à affronter.

La compétence, définie comme mobilisation de ressources internes et externes en situation, ne peut plus ignorer la question des outils cognitifs externalisés --- mais elle ne peut pas non plus s'y réduire: la compétence à l'ère des IAG est la capacité à interagir avec ces outils
de façon critique, réflexive et transférable. 
La réussite scolaire, partiellement construite sur une idéologie méritocratique déjà fragile, voit ses présupposés encore davantage exposés: le défi n'est pas de restaurer une méritocratie menacée, mais de se demander ce que l'école veut réellement certifier. 
L'évaluation, enfin, doit évoluer de la certification d'un produit fini vers l'observation d'un processus: non pas \og{}qu'est-ce que cet apprenant a produit?\fg{} mais \og{}comment cet apprenant construit-il sa pensée, avec les outils à sa disposition?\fg{}.

En perspective, l'enjeu dépasse la seule question des IAG. 
C'est la question plus ancienne du rôle de l'école face à l'évolution des outils cognitifs qui se pose ici avec une acuité nouvelle: de l'écriture à l'imprimerie, de la calculatrice à Internet, chaque outil a contraint l'école à redéfinir ce qu'elle entend cultiver. 
Les IA génératives ne sont probablement qu'une étape dans cette transformation longue. 
L'avenir de l'évaluation scolaire est peut-être moins dans la résistance à ces outils que dans l'invention de formes capables de saisir ce qu'ils ne peuvent pas faire à notre place: le jugement situé, la réflexivité sur son propre apprentissage, l'engagement éthique avec le savoir. 
Ce sont ces dimensions que l'école du XXI\textsuperscript{e}~siècle a plus que jamais vocation à cultiver.